% Figure 46.1 : où vit un Secret, et ce que chaque protection couvre (chapitre 46).
\documentclass[tikz,border=4pt]{standalone}
\usepackage{quai}
\begin{document}
\begin{tikzpicture}[x=1cm,y=1cm,
    b/.style={qbox, font=\scriptsize, align=left, inner sep=4pt, text width=#1, minimum height=1.05cm},
    p/.style={qbox, font=\scriptsize, align=left, inner sep=3pt, text width=#1, draw=QVert, fill=QVertBg},
    lien/.style={qlbl, font=\scriptsize, fill=QPaper, inner sep=1pt, align=center}]

  \node[b=3cm, draw=QBrique] (git) at (0,0) {\textbf{manifeste YAML}\\dans un dépôt Git :\\base64, donc en clair};
  \node[b=3.4cm, draw=QOcre, fill=QOcreBg] (api) at (5,0) {\textbf{API server}\\chiffre à l'écriture, déchiffre à la lecture (clé sur le nœud)};
  \node[b=3.2cm, draw=QBleu] (etcd) at (10,0) {\textbf{etcd}\\valeur actuelle, anciennes révisions};
  \node[b=3.2cm, draw=QBleu] (sav) at (14.6,0) {\textbf{sauvegardes}\\copies d'etcd, ailleurs};

  \node[b=3.4cm, draw=QViolet] (lect) at (5,-2.7) {\textbf{lecteurs de l'API}\\{\ttfamily get}, {\ttfamily list}, {\ttfamily watch} : en clair};
  \node[b=3.2cm, draw=QSarcelle] (kub) at (10,-2.7) {\textbf{kubelet}\\fichier en tmpfs, ou variable d'environnement};
  \node[b=3.2cm, draw=QSarcelle] (pod) at (14.6,-2.7) {\textbf{conteneur}\\lit la valeur en clair};

  \draw[qarr] (git) -- node[lien, above] {apply} (api);
  \draw[qarr] (api) -- (etcd);
  \draw[qarr] (etcd) -- (sav);
  \draw[qarr] (api) -- (lect);
  \draw[qarr] (api.south east) -- (kub.north west);
  \draw[qarr] (kub) -- (pod);

  \node[p=3cm] at (0,-4.7) {\textbf{Sealed Secrets}\\le dépôt ne contient que du chiffré};
  \node[p=3.4cm] at (5,-4.7) {\textbf{RBAC} (chapitre 43)\\qui peut lire par l'API};
  \node[p=3.2cm] at (10,-4.7) {\textbf{chiffrement au repos}\\etcd et ses sauvegardes, pas l'API};
  \node[p=3.2cm] at (14.6,-4.7) {\textbf{coffre externe}\\source de vérité hors du cluster, rotation};
\end{tikzpicture}
\end{document}
